\section{Exploratory Data Analysis}
\label{sec:proposed:eda}

The exploratory data analysis confirms what the schema in
Section~\ref{sec:proposed:schema} hints at: the AFC side is large,
clean, and analytically rich, while the GPS side is small (single
day, nine routes) and serves more as a proof-of-concept for the
GPS-driven modules until the full 50-day external dataset is
loaded. This section reports the dataset inventory, the AFC and
GPS findings that scope the dashboard design, and the data quality
issues that govern the constraints of
Section~\ref{sec:proposed:modules}.

\subsection{Dataset inventory}

Table~\ref{tab:proposed:dataset} summarises the data available
locally, confirmed against the production DuckDB instance on
2026-05-07.

\begin{table}[H]
\centering
\caption{Dataset inventory.}
\label{tab:proposed:dataset}
\begin{tabularx}{\linewidth}{lXr}
\hline
Stream & Period and coverage & Rows \\
\hline
AFC tickets        & 2026-01-01 to 2026-03-31 (90 days, 128 routes, 10 companies) & 17{,}883{,}623 \\
AFC payments       & same                                                          & 17{,}883{,}623 \\
GPS raw pings      & 2025-03-18 (single day, 131 vehicles)                         &    720{,}088 \\
GPS extracted trips& 2025-03-18, 9 routes (HMM-filtered)                           &      1{,}244 \\
GPS stop-times     & 2025-03-18                                                    &     12{,}345 \\
HCMC stops master  & static                                                        &      4{,}710 \\
\hline
\end{tabularx}
\end{table}

\subsection{AFC findings}

Five AFC-side findings shape the design.

\paragraph{Passenger-segment composition.}
The four demographic segments derived deterministically from
\texttt{ticket\_type} (Section~\ref{sec:proposed:schema}) are
distributed as shown in Table~\ref{tab:proposed:segments}. Three
observations stand out: the student segment alone accounts for
35.6\% of all boardings, free or subsidised passes account for
15.7\%, and the full-fare segment is essentially empty at 0.02\%.
The implication for the design is that the network operates on
government subsidy, and any analytical module that expects a
meaningful full-fare population on the local data is unsound.

\begin{table}[H]
\centering
\caption{Passenger segments (AFC, 2026 Q1).}
\label{tab:proposed:segments}
\begin{tabularx}{\linewidth}{Xrr}
\hline
Segment & Tickets & Share \\
\hline
Regular (subsidised smart-card)              & 8{,}708{,}026 & 48.7\% \\
Student (\texttt{HS-SV})                     & 6{,}362{,}862 & 35.6\% \\
Free / Subsidised (\texttt{miễn phí, CCCD})  & 2{,}809{,}726 & 15.7\% \\
Full fare (\texttt{không trợ giá})           &      3{,}009 & 0.02\% \\
\hline
\end{tabularx}
\end{table}

\begin{figure}[H]
\centering
\includegraphics[width=0.65\linewidth]{images/screenshots/eda_segment_donut}
\caption{Passenger-segment mix as a donut. The student slice is
the dominant single segment after subsidised regular riders;
the full-fare slice is barely visible.
Source: \texttt{notebooks/eda\_tickets.ipynb}.}
\label{fig:proposed:eda:segments}
\end{figure}

\paragraph{Payment channel.}
96.2\% of all tickets are recorded under the raw label
\texttt{`Tiền mặt'}. Field investigation confirmed that this label
denotes smart-card / prepaid deductions in this dataset, not
physical cash. Mastercard (1.5\%), Visa (1.1\%), elderly CCCD
(0.5\%), Momo (0.5\%) and ZaloPay (0.1\%) make up the remainder.
The label is renamed to \emph{Smart card / Prepaid} at the Silver
layer and consistently in every chart.

\paragraph{Revenue and fare distribution.}
Total revenue across the 90-day window is 72.45 billion VND,
with an average fare of 4{,}051 VND per trip. The fare
distribution is bimodal: a large mass at 3{,}000 VND
(subsidised rate) and a
secondary peak at the unsubsidised rate. Zero-price tickets (free
passes) account for 16.9\% of the row count.

\paragraph{Tri-peak temporal demand.}
Demand peaks at three times of day rather than the conventional
two: 06:00--07:00 (school commute), 11:00 (student midday return),
and 15:00--17:00 (school dismissal plus work return). The
top-six hours by total boardings are listed in
Table~\ref{tab:proposed:peakhours}. The 11:00 peak is
distinctively HCMC and follows directly from the high student
share of ridership.

\begin{table}[H]
\centering
\caption{Top six demand hours, 2026 Q1.}
\label{tab:proposed:peakhours}
\begin{tabularx}{0.8\linewidth}{Xrr}
\hline
Hour & Boardings & Network share \\
\hline
16:00 & 1{,}743{,}725 & 9.7\% \\
06:00 & 1{,}685{,}255 & 9.4\% \\
17:00 & 1{,}637{,}497 & 9.2\% \\
11:00 & 1{,}363{,}931 & 7.6\% \\
07:00 & 1{,}323{,}666 & 7.4\% \\
15:00 & 1{,}157{,}028 & 6.5\% \\
\hline
\end{tabularx}
\end{table}

\begin{figure}[H]
\centering
\includegraphics[width=0.95\linewidth]{images/screenshots/eda_hourly_demand}
\caption{Hourly boarding demand across the network, 2026 Q1.
Three peaks are visible at 06:00--07:00, 11:00, and
15:00--17:00 --- the distinctive HCMC tri-peak pattern.
Source: \texttt{notebooks/eda\_tickets.ipynb}.}
\label{fig:proposed:eda:hourly}
\end{figure}

\paragraph{Route concentration and ghost trips.}
The top-10 routes (out of 128) carry roughly 32\% of all
boardings, with routes 20 and 150 as clear outliers
(6.3\% and 5.5\% respectively). At the trip level, 213{,}408 of
1{,}199{,}766 unique trips (17.8\%) ran with fewer than five
boardings --- the near-ghost rate. Some routes exceed 40\%,
suggesting a structural mismatch between scheduled frequency and
actual demand on those routes.

\begin{figure}[H]
\centering
\includegraphics[width=0.95\linewidth]{images/screenshots/eda_ghost_distribution}
\caption{Near-ghost rate per route (sorted descending). The
network-wide rate of 17.8\% is exceeded by roughly a quarter
of routes, several of which run above 40\%.
Source: \texttt{notebooks/eda\_tickets.ipynb}.}
\label{fig:proposed:eda:ghost}
\end{figure}

\subsection{GPS findings}

The local GPS sample is a single day over nine routes. It is
sufficient to demonstrate the GPS-driven analytical modules but
not to compute their indicators at network scale. Three GPS
findings are relevant.

\paragraph{Raw pings.}
720{,}088 pings from 131 vehicles. The instantaneous speed
field is null for 49.7\% of pings: GPS devices do not emit a
speed value when the bus is stationary at a stop, in heavy
traffic, or parked. The non-null average is 23.9~km/h, which is
consistent with typical urban mixed-traffic conditions.

\paragraph{HMM extraction and QoS-2 caveat.}
The HMM-filtered \texttt{gps\_extracted} table covers 1{,}244
trips on 9 routes and contains 369{,}887 GPS points. Across this
set, the aggregate \emph{on-path rate} (the fraction of pings
inside the 50\,m route buffer used by QoS-2) is 25.1\%. This is
much lower than would be expected if the metric directly measured
service compliance: positioning noise in HCMC's dense urban
canyons causes pings to fall outside the buffer even when the bus
is in fact on-route. The dashboard surfaces this caveat as an
informational callout next to the QoS-2 KPI in the Route Explorer
lens; the metric is interpreted relatively (comparing routes to
each other) rather than absolutely.

\paragraph{Stop dwell.}
The \texttt{gps\_stop\_times} table records 12{,}345 stop visits
with median dwell of 30~seconds (consistent with a rapid-boarding
system). Outliers above 300~seconds may indicate mechanical issues,
high passenger volume, or schedule-holding behaviour at timing
points. Dwell statistics are available only for the 9 GPS routes
and the single sample day.

\subsection{Data-quality issues and constraints}

Five issues constrain the design.

\begin{table}[H]
\centering
\caption{Data-quality issues and the resulting design constraints.}
\label{tab:proposed:dq}
\renewcommand{\arraystretch}{1.25}
\begin{tabularx}{\linewidth}{p{4.4cm}X}
\hline
\textbf{Issue} & \textbf{Constraint imposed} \\
\hline
GPS (2025) and AFC (2026) cover different calendar windows.
& Treat as complementary sources. Combine narratively at the stop level rather than fuse at the trip level. \\

AFC records boarding events only; no alighting events.
& All crowding metrics framed as ``boardings per stop / per trip''. ``Passengers on the bus at time T'' is not a computable quantity. \\

GPS local sample = 1 day, 9 routes; full 50-day dataset is external.
& GPS-driven modules labelled ``(sample: 2025-03-18)''. GroupSet temporal clustering and holiday-vs-normal comparison wait for the full dataset. \\

QoS-2 absolute scores are systematically low (urban-canyon GPS noise).
& Surface as a relative metric only. Document the caveat next to the KPI. \\

\texttt{`Tiền mặt'} mislabel and 16.9\% zero-price tickets.
& Rename label at the Silver layer. Keep zero-price tickets as legitimate boardings (they are free passes, not data errors). \\
\hline
\end{tabularx}
\end{table}
